


\section{Introduction}

Large language models are increasingly deployed not as isolated text generators
but as components of larger \emph{AI agent} systems that perform structured
analytical tasks. In such systems, a language model operates within a broader
workflow that may include prompts, reasoning scaffolds, external tools,
verification procedures, and decision logic. These agent architectures are
being used for tasks such as policy analysis, incident investigation, software
diagnostics, and financial risk assessment.

Despite the rapid proliferation of AI agent architectures, methods for
evaluating them remain underdeveloped. Existing evaluation practices typically
fall into three categories. First, many studies rely on benchmark datasets
designed for question answering, reasoning puzzles, or code generation.
Although useful for measuring certain capabilities, such benchmarks do not
capture the open-ended analytical tasks that many AI agents are intended to
perform. Second, some work evaluates agent systems through qualitative case
studies or demonstrations, which can provide insight but lack experimental
control and reproducibility. Third, evaluations often compare outputs produced
under different prompts or agent designs without controlling for information
leakage, shared inputs, or evaluator bias.

These limitations make it difficult to determine whether a proposed agent
architecture genuinely improves analytical reasoning or whether observed
differences arise from uncontrolled variation in prompts, inputs, or evaluation
procedures.

This paper proposes a controlled experimental framework for evaluating AI
agents performing complex analytical tasks. The central idea is to isolate
agent interventions while holding all other aspects of the task environment
constant. The framework consists of several components:

\begin{itemize}

\item A frozen dataset of analytical tasks constructed from historical
socio-technical incidents across multiple domains.

\item Neutral task descriptions that describe system structure and task intent
without revealing investigative conclusions.

\item Controlled agent conditions in which different agent architectures are
applied to identical task inputs.

\item Independent reference materials derived from authoritative investigation
reports that serve as external analytical anchors.

\item A reproducible generation pipeline that produces structured analytical
artifacts under fixed model configurations.

\item Normalization and anonymization of outputs to remove stylistic cues and
enable blinded evaluation.

\item Pairwise human comparison of analytical outputs across multiple
evaluation dimensions.

\end{itemize}

The framework is instantiated using a dataset of twenty-five historical cases
including financial crises, aerospace accidents, technology failures, and
institutional breakdowns. These cases were selected because they possess rich
public documentation and independent investigation reports, enabling controlled
construction of analytical tasks.

Within this framework we compare three AI agent configurations. The first is a
baseline LLM agent that produces analysis directly from the task description.
The second is a protocol-guided agent that follows a structured reasoning
workflow. The third is an evidence-augmented agent that operates with excerpts
from independent investigation reports. Although the structured reasoning
protocol used in this study is one particular intervention, the experimental
framework is designed to evaluate a wide range of AI agent architectures.

The primary contribution of this work is methodological rather than
algorithmic. We present a reproducible experimental design for evaluating AI
agents on complex analytical tasks under controlled and blinded conditions.
This framework can be reused to evaluate prompt strategies, reasoning
protocols, tool-augmented workflows, and other forms of AI agent architecture.
\section{AI Agents and Interventions}

\subsection{AI Agents for Analytical Tasks}

In this work we use the term \emph{AI agent} to describe a structured
computational procedure that uses a language model as a reasoning component
within a broader analytical workflow. Rather than producing a single
response from a prompt, an AI agent may incorporate additional structure
such as prompts, reasoning procedures, tool interactions, or verification
steps.

Formally, an AI agent can be viewed as a mapping

\[
A : X \rightarrow Y
\]

where \(X\) represents a task input and \(Y\) represents a structured
analytical output. The mapping is implemented through a workflow that may
include intermediate artifacts, reasoning steps, or external information
sources.

Examples of AI agent structures include:

\begin{itemize}

\item Direct LLM analysis in which a language model generates an analytical
response from a task description.

\item Prompt scaffolds that guide the model through structured reasoning
steps.

\item Protocol-guided workflows that enforce ordering constraints on the
analysis process.

\item Tool-augmented agents that retrieve information or execute external
procedures during analysis.

\item Evidence-augmented agents that incorporate excerpts from external
documents or investigative materials.

\end{itemize}

In practice, modern AI systems frequently combine several of these elements,
forming a reasoning pipeline in which the language model acts as one
component of a larger analytical procedure.

\subsection{Interventions in Agent Workflows}

An \emph{intervention} is defined as a modification to the analytical
workflow of an AI agent that is intended to influence the quality,
structure, or reliability of its outputs.

Examples of interventions include:

\begin{itemize}

\item introducing structured reasoning prompts or protocols;

\item providing external evidence or retrieved documents;

\item enforcing ordering constraints on analytical steps;

\item requiring explicit articulation of assumptions or uncertainties;

\item adding verification or diagnostic stages to the workflow.

\end{itemize}

In experimental evaluation, the objective is to determine whether a given
intervention improves the analytical usefulness of the outputs produced by
the agent.

\subsection{Experimental Isolation of Interventions}

A central challenge in evaluating AI agents is isolating the effect of a
specific intervention. Differences in prompts, task descriptions, external
information, or evaluation procedures can easily confound experimental
results.

The framework proposed in this paper isolates interventions by holding the
following elements constant across experimental conditions:

\begin{itemize}

\item the task dataset;
\item the task input descriptions;
\item the underlying language model and generation configuration;
\item the evaluation procedure and scoring rubric.
\end{itemize}

Only the agent workflow itself is varied. This allows observed differences
in analytical outputs to be attributed to the intervention under study
rather than to unrelated experimental variation.

\subsection{Agent Configurations Evaluated in This Study}

To demonstrate the proposed experimental framework, we evaluate three agent
configurations that differ in their analytical workflow:

\begin{itemize}

\item \textbf{Baseline agent (B)}: the language model generates analysis
directly from the task description.

\item \textbf{Protocol-guided agent (P)}: the language model follows a
structured reasoning workflow that constrains the order and form of
analytical steps.

\item \textbf{Evidence-augmented agent (I)}: the language model performs
analysis using excerpts derived from independent investigation reports.

\end{itemize}

These configurations are used as an instantiation of the experimental
framework. The methodology itself is designed to evaluate a broad class of
AI agent architectures, including prompt scaffolds, reasoning protocols,
tool-augmented agents, and other structured workflows.
\section{Experimental Framework}

This section describes the experimental framework used to evaluate AI agent
workflows under controlled conditions. The framework is designed to isolate
the effect of agent interventions while ensuring reproducibility,
comparability of outputs, and independence of evaluation.

The experimental design consists of the following components:

\begin{enumerate}

\item construction of a frozen task dataset;

\item creation of neutral task inputs that describe the system under
analysis without revealing investigative conclusions;

\item definition of multiple agent configurations that differ only in
their analytical workflow;

\item generation of analytical artifacts using a controlled model
configuration;

\item normalization and anonymization of outputs;

\item blinded human evaluation using pairwise comparison;

\item statistical analysis of comparative preferences.

\end{enumerate}

Together these elements provide a reproducible procedure for experimentally
evaluating AI agents performing complex analytical tasks.

\subsection{Frozen Task Dataset}

The experimental task set consists of twenty-five historical cases drawn
from multiple socio-technical domains, including aerospace accidents,
financial failures, energy infrastructure incidents, technology failures,
and institutional governance breakdowns.

Cases were selected through an iterative discovery and screening process
designed to produce a diverse and information-rich evaluation set. Each
candidate case was required to satisfy the following criteria:

\begin{itemize}

\item the event involved a major failure or systemic breakdown;

\item sufficient public documentation existed to describe the system
structure and operational context;

\item at least one credible independent investigation or assessment report
was available;

\item the case was suitable for comparative analytical reconstruction.

\end{itemize}

The final set of twenty-five cases was assigned stable identifiers
(\texttt{C01}--\texttt{C25}) and frozen in the file
\texttt{case\_list\_v1.txt}. This file serves as the root index for all
subsequent experimental artifacts.

\subsection{Neutral Task Inputs}

For each case, a neutral description of the system and task context was
constructed. These descriptions focus on system architecture, operational
processes, and organizational structure while intentionally excluding
investigative conclusions or causal findings.

The construction procedure consisted of:

\begin{enumerate}

\item writing a short system description (target length $\leq$150 words);

\item removing statements that reveal the outcome of the investigation;

\item excluding conclusions drawn by investigative reports;

\item preserving only information necessary to understand the system
structure and task intent.

\end{enumerate}

The resulting descriptions were frozen in the artifact
\texttt{neutral\_case\_descriptions\_v1.txt}. During generation, these
descriptions are converted into individual task input files
(\texttt{Cxx\_case\_input.txt}) corresponding to each case.

\subsection{Independent Reference Materials}

To provide an external analytical anchor, authoritative investigation
sources were identified for each case. Sources were selected according to a
credibility hierarchy that prioritizes formal investigative authority:

\begin{enumerate}

\item court findings or judicial determinations;

\item congressional or parliamentary investigation reports;

\item regulatory authority reports;

\item independent investigation boards or commissions.

\end{enumerate}

For each case, the most authoritative available source was designated as
\emph{Primary A}. When a second broad investigation source was available,
an alternate source was designated as \emph{Primary B} for sensitivity
analysis. Supplemental sources were included when they provided additional
technical or operational detail not covered in the primary report.

Relevant excerpts describing system architecture, operational processes,
decision structures, causal mechanisms, and investigative findings were
manually extracted from these sources and stored as frozen artifacts
(\texttt{Cxx\_refA\_extract.txt}, \texttt{Cxx\_refB\_extract.txt}).

\subsection{Agent Configurations}

Three AI agent configurations are evaluated within the framework:

\begin{itemize}

\item \textbf{Baseline agent (B)}: the language model generates analysis
directly from the neutral task description.

\item \textbf{Protocol-guided agent (P)}: the model follows a structured
reasoning workflow that constrains the order and form of analytical
artifacts.

\item \textbf{Evidence-augmented agent (I)}: the model performs analysis
using excerpts from independent investigation materials.

\end{itemize}

All agent configurations operate on identical task inputs and use the same
underlying model configuration. The only difference between conditions is
the analytical workflow applied by the agent.

\subsection{Reproducible Generation Pipeline}

Analytical outputs are generated using a controlled execution pipeline
implemented in a reproducible notebook environment.

For each case identifier \(C_{xx}\), the pipeline constructs a task input
file:

\[
C_{xx}\_\text{case\_input.txt}
\]

Generation is then performed for each agent configuration, producing raw
analytical artifacts:

\begin{itemize}

\item \texttt{Cxx\_B\_raw.txt}

\item \texttt{Cxx\_P\_raw.txt}

\item \texttt{Cxx\_I\_refA\_raw.txt}

\item \texttt{Cxx\_I\_refB\_raw.txt}

\end{itemize}

All generations are executed under a fixed model configuration specified in
\texttt{model\_config\_v1.json}. Each generated output is stored together
with a SHA-256 hash and a generation log entry recording the case
identifier, generation condition, model configuration, and timestamp.

\subsection{Output Normalization and Blinding}

Raw outputs produced by different agents may vary in length, structure, and
stylistic presentation. To enable fair comparison, outputs are normalized
into a common analytical template and anonymized prior to evaluation.

Each normalized document is assigned a random identifier (e.g.,
\texttt{D0001}, \texttt{D0002}) that hides the generation condition from
human evaluators. The mapping between document identifiers and generation
conditions is stored in a separate key file not accessible to raters.

\subsection{Blinded Human Evaluation}

Human raters evaluate outputs using pairwise comparison. Rather than
scoring individual documents in isolation, raters are presented with two
analytical outputs describing the same case and asked to indicate which
analysis is closer to a high-quality analytical reconstruction.

Each comparison is evaluated along multiple dimensions, including:

\begin{itemize}

\item analytical clarity;

\item structural reasoning quality;

\item diagnostic usefulness;

\item overall analytical usefulness.

\end{itemize}

Document order is randomized to prevent systematic position bias.

\subsection{Statistical Analysis}

Evaluation results are analyzed using pairwise preference statistics.
Directional preferences between agent conditions are aggregated across
cases and raters.

The primary measure is the preference rate between conditions. Statistical
significance of observed preference differences is assessed using
non-parametric tests such as the binomial sign test.

This analysis allows the experimental framework to quantify whether a given
agent intervention leads to systematically preferred analytical outputs.
\section{Dataset Construction}

The experimental framework requires a task dataset that supports controlled
analytical reconstruction while remaining independent of the evaluation
procedures applied by the agents. To satisfy this requirement, a dataset of
twenty-five historical socio-technical cases was constructed and frozen
prior to generation and evaluation.

The objective of the dataset construction process was not to create a
statistically representative sample of failures, but rather to assemble an
information-rich collection of cases that support structured analysis and
independent verification.

\subsection{Case Discovery}

Candidate cases were identified through an iterative discovery process
covering multiple socio-technical domains, including:

\begin{itemize}

\item finance and financial regulation;

\item aerospace and spaceflight;

\item energy infrastructure;

\item transportation systems;

\item technology and software systems;

\item corporate governance and institutional systems.

\end{itemize}

The initial discovery stage prioritized breadth of coverage rather than
immediate inclusion. Candidate events were drawn from publicly documented
historical failures and controversies known to involve complex system
interactions, organizational decision-making, or governance breakdowns.

\subsection{Shortlisting Criteria}

Candidate cases were then screened against four requirements.

\begin{enumerate}

\item \textbf{Major system failure or breakdown}

The event must involve substantial consequences such as financial loss,
loss of life, mission failure, systemic disruption, or significant
governance failure.

\item \textbf{Publicly describable system setup}

Sufficient public documentation must exist to describe the system
architecture, operational processes, and organizational structure without
relying on the conclusions of investigative reports.

\item \textbf{Availability of independent assessment}

The case must have at least one credible external investigation or
assessment source, such as court findings, congressional investigations,
regulatory reports, or independent inquiry commissions.

\item \textbf{Suitability for analytical reconstruction}

The case must support reconstruction of system behavior and decision
processes in a way that allows meaningful comparison between different
analytical workflows.

\end{enumerate}

\subsection{Coverage of the Solution Space}

During shortlisting, cases were intentionally selected to span a wide
solution space rather than cluster around a single failure type.

The final dataset includes cases characterized by:

\begin{itemize}

\item mechanistic engineering failures;

\item socio-technical governance failures;

\item financial and institutional breakdowns;

\item regulatory and oversight failures;

\item software and information-system failures.

\end{itemize}

This diversity allows the experimental framework to evaluate whether
different AI agent workflows perform differently across classes of
analytical problems.

\subsection{Exclusion Criteria}

Candidate cases were excluded if they failed to meet one or more of the
following conditions:

\begin{itemize}

\item no credible independent investigation source could be identified;

\item publicly available documentation of system structure was too limited;

\item the case was redundant with another selected case;

\item the case remained too unresolved or factually unstable to support a
frozen analytical task.

\end{itemize}

\subsection{Case Identification and Freezing}

After iterative screening, twenty-five cases were selected and assigned
stable identifiers:

\[
C01, C02, \ldots, C25
\]

The finalized case list was frozen in the artifact

\texttt{case\_list\_v1.txt}

which serves as the root index for all experimental inputs and generated
artifacts.

Each case identifier is subsequently used to generate associated task
inputs, reference materials, and analytical outputs throughout the
experimental pipeline.

\subsection{Neutral Case Descriptions}

For each case, a neutral system description was created to serve as the
primary task input for agent analysis.

The construction procedure consisted of:

\begin{enumerate}

\item writing a concise description of the system and operational context
(target length $\leq$150 words);

\item focusing on system architecture, operational processes, and task
intent;

\item removing investigative conclusions and causal findings;

\item excluding language that directly reveals the outcome of the event.

\end{enumerate}

These descriptions were frozen in the artifact

\texttt{neutral\_case\_descriptions\_v1.txt}

and later converted into individual case input files during generation.

The purpose of this neutralization procedure is to ensure that all agent
configurations operate from the same initial information state, preventing
direct leakage of investigative conclusions into the analytical task.

\subsection{Dataset Freezing and Reproducibility}

To ensure reproducibility, all dataset artifacts were frozen prior to model
generation and evaluation. Each artifact was recorded together with a
cryptographic hash in a freeze manifest.

The frozen artifacts include:

\begin{itemize}

\item \texttt{case\_list\_v1.txt}

\item \texttt{neutral\_case\_descriptions\_v1.txt}

\item \texttt{prompt\_baseline\_v1.txt}

\item \texttt{prompt\_protocol\_v1.txt}

\item \texttt{prompt\_independent\_v1.txt}

\item \texttt{protocol\_specification\_v1.txt}

\item \texttt{model\_config\_v1.json}

\item \texttt{case\_id\_mapping\_v1.csv}

\end{itemize}

Hashes for these artifacts are recorded in

\texttt{case\_freeze\_manifest\_v1.txt}

which enables verification that the dataset and experimental configuration
remain unchanged across experimental runs.
